\section{Extensions to embedding dimension greater than four}
\label{sec:extensions}

The Ap\'ery framework extends to all embedding dimensions. The geometric
surplus argument does not presently do so. We separate these two facts:
this section gives identities and a conditional transfer principle, but
does not claim Wilf's conjecture for any unrestricted \(e>4\).

\subsection{Dimension-independent moment and collision identities}

Let \(e=d+1\), \(S=\langle m,a_1,\ldots,a_d\rangle\), and define
\(A,w,H,T,s\) as before, with \(T\subseteq\mathbb N^d\). Preferred
factorizations again give \(|T|=m\), lower closure, and distinct residues.
For the indexed family of coordinate lines, the same progression sums give
\[
\sum_\ell L_\ell=dm,\qquad
\sum_\ell L_\ell t_\ell=(d+1)s.
\]
Consequently Zhai's weighted downset baseline \cite{Z} reads
\[
D_e:=dmH-(d+1)w\cdot s\ge0.
\]
The genus formula and \(W_e=en-c=(e-1)c-eg\) give the exact identity
\begin{equation}\label{eq:extension-moment}
mW_e=A D_e-\frac{e-2}{2}\,m(m-1).
\end{equation}
The factor \(m(m-1)\) makes the correction integral even when its
coefficient is a half-integer.

The predecessor-subtraction collision proof also remains valid. A minimal
excluded point has support disjoint from that of its representative; two
distinct minimal exclusions with intersecting support have distinct
residues. In particular there is still at most one full-support corner,
with representative zero. Unlike in three exponent coordinates, the
remaining corners can have support three or more.

Finally, if nonnegative masses on \(T\) have total \(dm\) and moment
\((d+1)s+\beta\), with \(\beta\ge0\), their destinations have weight
at most \(H\), and \(w_i\ge1\), then
\[
D_e=\sum_x\alpha_x(H-w\cdot x)+w\cdot\beta\ge|\beta|_1.
\]
Thus the local centroid mechanism extends. What is missing is a uniform
surplus of the required size for every relevant ideal.

\subsection{A conditional mesh-transfer principle}

The translated-mesh argument has a useful dimension-independent form.
Let \(K\) be a measurable positive-volume downset in the unit simplex.
Choose an integer \(N>d\), put \(h=1/N\), and define
\(T_z=\{y\in\mathbb N^d:z+hy\in K\}\).
Suppose that, for almost every \(z\in[0,h)^d\), this induced discrete
ideal satisfies
\[
d r_z|T_z|-(d+1)\sum_{y\in T_z}|y|_1\ge\alpha|T_z|,
\qquad
r_z=\left\lfloor\frac{1-|z|_1}{h}\right\rfloor.
\]
Only allowances \(N-d,\ldots,N-1\) occur almost everywhere. Writing
\(\epsilon_z=1-|z|_1-hr_z\), the integrand expands as
\[
\begin{aligned}
\sum_{y\in T_z}\bigl(d-(d+1)|z+hy|_1\bigr)
&=hD_{r_z}(T_z)+(d\epsilon_z-|z|_1)|T_z|\\
&\ge(\alpha h-|z|_1)|T_z|.
\end{aligned}
\]
After integrating, each coordinate remainder has mean at most \(h/2\)
by the interval-fiber calculation in Section~\ref{sec:p11}. Hence
\begin{equation}\label{eq:extension-mesh}
d-(d+1)\mathbb E_K|X|_1
\ge\left(\alpha-\frac d2\right)\frac1N.
\end{equation}
Our strip transfer is \(d=3\), \(N=21\), \(\alpha=4\), giving \(5/42\).
In higher dimension this is a conditional tool: the required finite
inequality and an exhaustive way to check it have not been established
here. A positive gap needs \(\alpha>d/2\).

\subsection{Why the three-horn decomposition does not extend directly}

The compatibility graph is chordal in three coordinate classes because
an induced cycle of length at least four repeats a class. With four
classes, that argument fails, and the failure is real.

Consider the finite lower ideal in \(\mathbb N^4\)
\[
U=\{x\in\{0,1\}^4:x_1x_3=0,\ x_2x_4=0\}.
\]
Its minimal exclusions are \(2e_i\) for \(1\le i\le4\), and
\(e_1+e_3,e_2+e_4\). None has full support. The compatibility graph
on the four positive axis levels has the induced cycle
\[
(1,1)\;\text{---}\;(2,1)\;\text{---}\;(3,1)\;\text{---}\;(4,1)\;\text{---}\;(1,1),
\]
whose two diagonals are absent.

More directly, no central box \([0,c]\), \(c\in U\), together with
four horns restricted to exceeding only one central coordinate can
cover \(U\). The support of \(c\) is empty, a singleton, or an adjacent
pair of this cycle. In every case there is an adjacent pair disjoint
from that support. Its vector \(e_i+e_j\) lies in \(U\), exceeds two
central caps, and belongs to neither the central box nor any such horn.

This is an obstruction for arbitrary lower ideals, not an assertion
that \(U\) admits an Ap\'ery labeling. A higher-dimensional proof might
exploit additional arithmetic restrictions or replace the decomposition
with a richer structural model. The current three-horn recurrence cannot
simply be reused with more axes.

Known dimension-independent results still cover substantial regions,
including \(3e\ge m\) \cite{Graph}, \(c\le3m\) \cite{Macaulay}, and
\(t\le e-1\), where \(t\) is the type \cite{FGH}. They do not supply
the missing uniform geometric surplus for a fixed \(e>4\).
